\documentclass[aps,preprint]{revtex4}%
\usepackage{amssymb}
\usepackage{amsfonts}
\usepackage{amsmath}
\usepackage{graphicx}%
\setcounter{MaxMatrixCols}{30}
%TCIDATA{OutputFilter=latex2.dll}
%TCIDATA{Version=5.50.0.2960}
%TCIDATA{CSTFile=revtex4.cst}
%TCIDATA{Created=Monday, January 18, 2016 14:25:21}
%TCIDATA{LastRevised=Tuesday, September 24, 2024 06:38:30}
%TCIDATA{<META NAME="GraphicsSave" CONTENT="32">}
%TCIDATA{<META NAME="SaveForMode" CONTENT="1">}
%TCIDATA{BibliographyScheme=Manual}
%TCIDATA{<META NAME="DocumentShell" CONTENT="Articles\SW\REVTeX 4">}
%BeginMSIPreambleData
\providecommand{\U}[1]{\protect\rule{.1in}{.1in}}
%EndMSIPreambleData
\newtheorem{theorem}{Theorem}
\newtheorem{acknowledgement}[theorem]{Acknowledgement}
\newtheorem{algorithm}[theorem]{Algorithm}
\newtheorem{axiom}[theorem]{Axiom}
\newtheorem{claim}[theorem]{Claim}
\newtheorem{conclusion}[theorem]{Conclusion}
\newtheorem{condition}[theorem]{Condition}
\newtheorem{conjecture}[theorem]{Conjecture}
\newtheorem{corollary}[theorem]{Corollary}
\newtheorem{criterion}[theorem]{Criterion}
\newtheorem{definition}[theorem]{Definition}
\newtheorem{example}[theorem]{Example}
\newtheorem{exercise}[theorem]{Exercise}
\newtheorem{lemma}[theorem]{Lemma}
\newtheorem{notation}[theorem]{Notation}
\newtheorem{problem}[theorem]{Problem}
\newtheorem{proposition}[theorem]{Proposition}
\newtheorem{remark}[theorem]{Remark}
\newtheorem{solution}[theorem]{Solution}
\newtheorem{summary}[theorem]{Summary}
\newenvironment{proof}[1][Proof]{\noindent\textbf{#1.} }{\ \rule{0.5em}{0.5em}}
\begin{document}
\preprint{HEP/123-qed}
\title[Short title for running header]{REV\TeX4}
\author{First Author}
\affiliation{My Institution}
\author{Second Author}
\affiliation{My Institution}
\author{Third Author}
\affiliation{Other Institution}
\keywords{one two three}
\pacs{PACS number}

\begin{abstract}
Shell document for REV\TeX{} 4.

\end{abstract}
\volumeyear{year}
\volumenumber{number}
\issuenumber{number}
\eid{identifier}
\date[Date text]{date}
\received[Received text]{date}

\revised[Revised text]{date}

\accepted[Accepted text]{date}

\published[Published text]{date}

\startpage{101}
\endpage{102}
\maketitle
\tableofcontents


\section{Types of Von Neumann Algebras}

\begin{remark}
Algebraically a von Neumann algebra is a $\ast$-algebra of bounded operators
on some Hilbert space of countable dimension and is the commutant of its commutant.
\end{remark}

\begin{remark}
Topologically a von Neumann algebra is an $\ast$-algebra of operators that is
closed in the \emph{weak operator topology. }A sequence of operators $\left\{
T_{i}\right\}  $ converges in the \emph{weak operator topology (WOT)} to an
operator $T$ if their expectation values converge to that of $T$ in every
state, that is,
\begin{equation}
\lim\limits_{i\rightarrow\infty}\left\langle \left\vert \Psi\right\vert
T_{i}\Psi\right\rangle =\left\langle \left\vert \Psi\right\vert T\Psi
\right\rangle \forall\,\Psi\in\mathcal{H}.
\end{equation}

\end{remark}

\emph{It's a nontrivial theorem that these two definitions agree!}

Classifying all $\ast$-algebras of operators is an utterly hopeless task,
while classifying von Neumann algebras is almost within reach. Every von
Neumann algebra can be built from so-called \emph{simple} ones as a direct
sum, or more generally a \emph{direct integral}. As usual in algebra, the
simple von Neumann algebras are defined to be those without any nontrivial
ideals. This turns out to be equivalent to saying that only scalar multiples
of the identity commute with everything in the von Neumann algebra.

Simple von Neumann algebras are called \emph{factors}. Thus we just need to
classify the factors: the process of combining these together to get the other
von Neumann algebras is well defined.

The first step in classifying factors was done by von Neumann and Murray, who
divided them into Types $I,II,$and $III$. This classification involves the
concept of a trace or weight, which is a generalization of the usual trace of
a matrix.

\begin{definition}
A (faithful, normal) trace on $\mathcal{A}$ is a function%
\begin{equation}
\U{3c4} :\mathcal{A\rightarrow}\mathbb{C}%
\end{equation}
satisfying:

\textbullet\qquad$\U{3c4} $ is a linear map;

\textbullet\qquad$\U{3c4} $ is positive, that is,
\begin{equation}
\U{3c4} (x^{\dagger}x)\geq0\forall x\in\mathcal{A}%
\end{equation}
i.e
\begin{equation}
\tau:\mathcal{C}\rightarrow\left[  0,\infty\right]  ;
\end{equation}


\textbullet\qquad(normality) $\U{3c4} $ is weakly continuous (the identity is
its corresponding density operator);

\textbullet\qquad(faithful) $\U{3c4} (x^{\dagger}x)$ $=0$ if and only if $x=0$;

\textbullet\qquad(trace property) $\U{3c4} (xy)=\U{3c4} (yx)$.
\end{definition}

A trace is very much like an integral, so we're really studying a
noncommutative version of the theory of integration. On the other hand, in the
matrix case, the trace of a projection operator is just the dimension of the
space it's the projection onto. We can define a "projection" in any von
Neumann algebra, $\mathcal{A},$ to be an operator with
\begin{equation}
p=p^{\dagger}=p^{2}.
\end{equation}
Traces, weights and states can be generalized to \emph{centre valued} linear
functions i.e.
\begin{equation}
\U{3c4} :\mathcal{A\rightarrow Z}\left(  \mathcal{A}\right)  .
\end{equation}


\begin{definition}
For $\mathcal{A}\subset\mathcal{B(H)}$ a von Neumann algebra, a
\emph{center-valued state} is a linear map $\U{3d5} :\mathcal{A\rightarrow
Z}\left(  \mathcal{A}\right)  $ such that:
\end{definition}

\begin{enumerate}
\item $\U{3d5} (x^{\dagger}x)\geq0$ for all $x\in\mathcal{A}$

\item $\U{3d5} (z)=z$ for all $z\in\mathcal{Z}\left(  \mathcal{A}\right)  $

\item $\U{3d5} (zx)=z\U{3d5} (x)$ for all $x\in\mathcal{A}$ and $z\in
\mathcal{Z}\left(  \mathcal{A}\right)  $

\item $\U{3d5} $ is \emph{normal} if it is $\U{3c3} -$WOT continuous (ultra
weak continuous), normal can also be defined as states that preserves suprema
when applied to increasing nets of self adjoint operators; or equivalently to
increasing sequences of projections.
\end{enumerate}

Scalar valued states are special cases of center values states.

\paragraph{Weights, states, and traces}

Weights and their special cases states and traces are discussed in detail in
(Takesaki 1979).

A \emph{weight} $\U{3c9} $ on a von Neumann algebra is a linear map from the
set of positive elements (those of the form $x^{\dagger}x$) to $[0,\infty]$.

A positive linear functional is a weight with $\U{3c9} (I_{\mathcal{A}})$
finite (or rather the extension of $\U{3c9} $ to the whole algebra by linearity).

\begin{enumerate}
\item A \emph{state} is a weight with $\U{3c9} (I_{\mathcal{A}}%
)=I_{\mathcal{A}}$.

\item A \emph{trace} is a weight with $\U{3c9} (x^{\dagger}x)=\U{3c9}
(xx^{\dagger})$ for all $x$.

\item A \emph{tracial state} is a trace with $\U{3c9} (I_{\mathcal{A}%
})=I_{\mathcal{A}}$.
\end{enumerate}

Any factor has a trace such that the trace of a non-zero projection is
non-zero and the trace of a projection is infinite if and only if the
projection is infinite. Such a trace is unique up to rescaling. For factors
that are separable or finite, two projections are equivalent if and only if
they have the same trace. The type of a factor can be read off from the
possible values of this trace on projections as follows:

\begin{enumerate}
\item Type $I_{n}$: $0,x,2x,....,nx$ for some positive $x$ (usually normalized
to be $1/n$ or $1$); $n\in\mathbb{N}.$

\item Type $II_{1}:[0,x]$ for some positive $x$ (usually normalized to be $1$).

\item Type$II_{\infty}:[0,\infty]$.

\item Type $III$: $\left\{  0,\infty\right\}  $.
\end{enumerate}

\begin{theorem}
Every von Neumann algebra $\mathcal{A}\subset\mathcal{B(H)}$ has a normal,
center-valued state.
\end{theorem}

The set of all projections in a von Neumann algebra, $\mathcal{A},$ is denoted
$\mathcal{P}\left(  \mathcal{A}\right)  .$ If we study the trace of
projections, we're studying a GENERALIZATION OF THE CONCEPT OF DIMENSION. It
turns out this can be infinite, or even nonintegral.

Factors are:

\begin{enumerate}
\item Type $I_{n}$ if
\begin{equation}
\tau:\mathcal{P}\left(  \mathcal{A}\right)  \rightarrow\left\{  0,n\right\}
\subset\mathbb{N};1\leq n\leq\infty.
\end{equation}
It turn out that every Type $I_{n}$ factor is isomorphic to the algebra of
$n\times n$ matrices. Also, every type Type $I_{\infty}$ factor is isomorphic
to the algebra of all bounded operators on a Hilbert space of countably
infinite dimension.

\item Type $II_{1}$ if
\begin{equation}
\tau:\mathcal{P}\left(  \mathcal{A}\right)  \rightarrow\left[  0,1\right]
\subset\mathbb{R}.
\end{equation}
Type $II_{\infty}$ if
\begin{equation}
\tau:\mathcal{P}\left(  \mathcal{A}\right)  \rightarrow\left[  0,\infty
\right]  \subset\mathbb{R}.
\end{equation}
Playing with Type $II$ factors amounts to letting dimension be a continuous
rather than discrete parameter! It is easy to construct a Type $II_{1}$
factor
\begin{align}
\tau &  :%
%TCIMACRO{\tbigcup \limits_{1\leq n\leq r}}%
%BeginExpansion
{\textstyle\bigcup\limits_{1\leq n\leq r}}
%EndExpansion
M\left(  2^{n},\mathbb{C}\right)  \rightarrow\mathbb{C}\\
\left.  \tau\right.  _{n}  &  =\frac{1}{2^{n}}\left.  \operatorname{Tr}%
\right.  _{n}%
\end{align}
then
\begin{equation}
\mathcal{A}_{II_{1}}=\lim\limits_{n\rightarrow\infty}%
%TCIMACRO{\tbigcup \limits_{1\leq n\leq r}}%
%BeginExpansion
{\textstyle\bigcup\limits_{1\leq n\leq r}}
%EndExpansion
M\left(  2^{n},\mathbb{C}\right)  .
\end{equation}
One can show this von Neumann algebra is a factor. This isn't the only Type
$II_{1}$ factor, but it's the only one that contains a sequence of
finite-dimensional von Neumann algebras whose union is dense in the weak
topology. A von Neumann algebra like that is called "\emph{hyperfinite}", in
this case is called "\emph{the hyperfinite Type} $II_{1}$ \emph{factor}".

\begin{enumerate}
\item The algebra of $2^{n}\times2^{n}$ matrices is a Clifford algebra, so the
hyperfinite $II_{1}$ factor is a kind of infinite-dimensional Clifford algebra.

\item The Clifford algebra of $2^{n}\times2^{n}$ matrices is just another name
for the algebra generated by creation and annihilation operators on the
fermionic Fock space over $\mathbb{C}^{n}$. One can show that the hyperfinite
$II_{1}$ factor is the smallest von Neumann algebra containing the creation
and annihilation operators on a fermionic Fock space of countably infinite dimension.

\item There is more than one Type $II_{\infty}$ factor, but again there is
only one that is hyperfinite. You can get this by tensoring the Type
$I_{\infty}$ and the hyperfinite $II_{1}$.
\end{enumerate}

\item Type $III$ if all projections other than $0$ have infinite trace. In
other words. The only trace on type III factors takes the value $\infty$ on
all non-zero positive elements, and any two non-zero projections are equivalent.
\end{enumerate}

\begin{definition}
Lattice of Projections%
\begin{align}%
%TCIMACRO{\tbigwedge \limits_{i\in I}}%
%BeginExpansion
{\textstyle\bigwedge\limits_{i\in I}}
%EndExpansion
p_{i}  &  =\text{orthogonal projection onto }%
%TCIMACRO{\tbigcap \limits_{i\in I}}%
%BeginExpansion
{\textstyle\bigcap\limits_{i\in I}}
%EndExpansion
\text{ }p_{i}\mathcal{H}\\%
%TCIMACRO{\tbigvee \limits_{i\in I}}%
%BeginExpansion
{\textstyle\bigvee\limits_{i\in I}}
%EndExpansion
p_{i}  &  =\text{orthogonal projection onto }%
%TCIMACRO{\tbigcup \limits_{i\in I}}%
%BeginExpansion
{\textstyle\bigcup\limits_{i\in I}}
%EndExpansion
\text{ }p_{i}\mathcal{H}%
\end{align}

\end{definition}

\begin{definition}
For $\mathcal{A}\subset\mathcal{B(H)}$ a von Neumann algebra, a projection
$p\in\mathcal{P(A)}$ is said to be monic if there exists a finite collection
of pairwise orthogonal projections $\{p_{1},...,p_{n}\}\subset\mathcal{P(A)}$
such that $p_{i}\sim p$ for $1\leq i\leq n$ and $%
%TCIMACRO{\tsum \limits_{1\leq i\leq n}}%
%BeginExpansion
{\textstyle\sum\limits_{1\leq i\leq n}}
%EndExpansion
p_{i}\in\mathcal{Z}\left(  \mathcal{A}\right)  .$
\end{definition}

\begin{definition}
Let $\mathcal{A}\subset\mathcal{B}(\mathcal{H})$ be a von Neumann algebra,
then $p\in\mathcal{P}(\mathcal{A})$ is said to be
\end{definition}

\textbullet\qquad\emph{minimal} if $p\neq0$ and $q\leq p$ implies $q=0$ or
$q=p$; equivalently, $\dim\left\{  p\mathcal{A}p\right\}  =1$.

\textbullet\qquad\emph{abelian} if $p\mathcal{A}p$ is abelian.

\textbullet\qquad\emph{finite} if $q\leq p$ and $q\sim p$ implies $p=q$.

\textbullet\qquad\emph{semi-finite} if there exists a family $\{\left.
p_{i}\right\vert i\in I\}$ of pairwise orthogonal, finite projections such
that
\begin{equation}
p=%
%TCIMACRO{\tsum \limits_{i\in I}}%
%BeginExpansion
{\textstyle\sum\limits_{i\in I}}
%EndExpansion
p_{i}%
\end{equation}


\textbullet\qquad\emph{purely infinite} if $p\neq0$ and there does not exist
any non-zero finite projections $q\leq p$.

\textbullet\qquad\emph{properly infinite} if $p\neq0$ and for all non-zero
$z\in\mathcal{P}(\mathcal{Z}(\mathcal{A}))$ the projection $zp$ is not finite.

For a projection $p\in\mathcal{P}(\mathcal{A})$ we have the following
implications:%
\begin{equation}
\text{minimal}\Rightarrow\text{abelian}\Rightarrow\text{finite}\Rightarrow
\text{semi-finite}\Rightarrow\text{not purely infinite}%
\end{equation}


and%
\begin{equation}
\text{purely infinite}\Rightarrow\text{properly infinite.}%
\end{equation}


\begin{definition}
For $p,q\in\mathcal{P}(\mathcal{A})$, we say that p is sub-equivalent to $q$
(in $\mathcal{A}$ ) and write $p\preccurlyeq q$ if there exists a partial
isometry $v\in M$ such that $v^{\dagger}v=p$ and $vv^{\dagger}\leq q$. We say
$p$ is equivalent to $q$ (in $\mathcal{A}$ ) and write $p\sim q$ if there
exists a partial isometry $v\in\mathcal{A}$ such that $v^{\dagger}v=p$ and
$vv^{\dagger}\leq q$. If $p\preceq q$ but $p\nsim q$, we write $p\prec q$.
\end{definition}

\begin{remark}
If $p,q\in\mathcal{P}(\mathcal{A})$ are such that $p\leq q$, then $p\sim q$
with partial isometry $v=p$.
\end{remark}

\begin{proposition}
For $\mathcal{A}\subset\mathcal{B}(\mathcal{H})$, the relation $\sim$ is a
partial ordering and an equivalence relation on $\mathcal{P}(\mathcal{A})$.
\end{proposition}

\begin{definition}
A von Neumann algebra $\mathcal{A}\subset\mathcal{B}(\mathcal{H})$ is said to be
\end{definition}

\textbullet\qquad Type $I$ if every non-zero projection has a non-zero abelian subprojection.

\textbullet\qquad Type $II$ if it is semi-finite and has no non-zero abelian projections.

\textbullet\qquad Type $III$ if it is purely infinite.

\begin{theorem}
Let $\mathcal{A}\subset\mathcal{B}(\mathcal{H})$ be a von Neumann algebra.
Then there exists unique, central, pairwise orthogonal projections
$\mathcal{Z}I,\mathcal{Z}II,\mathcal{Z}III\in\mathcal{Z}(\mathcal{A})$ such
that $\mathcal{Z}I+\mathcal{Z}II+\mathcal{Z}III$ $=I$ and $\mathcal{Z}T$ is
Type $T$ for $T\in\{I,II,III\}$.
\end{theorem}


\end{document}